\section{The claim and its setting}\label{sec:introduction}

Resolution over parities augments ordinary resolution with affine equations
over \(\F\). The system was introduced by Itsykson and Sokolov
\cite{IS14,IS20}. Its unrestricted proof-size lower-bound problem remains an
explicit open benchmark in the literature reviewed for this draft
\cite{IPS26}. The contribution claimed here is a superpolynomial lower
bound for the standard bit pigeonhole principle in unrestricted DAG-like
\(\Res\).

\subsection{The bit pigeonhole principle}
Fix \(\ell\ge2\), put \(n=2^\ell\) and \(m=n+1\), and introduce
\[
  b_{it}\qquad (i\in[m],\ t\in[\ell]).
\]
The row \(b_i\) encodes one of the \(n\) labels in \(\{0,1\}^{\ell}\).
For a label \(z\), let \([b_{it}\ne z_t]\) denote \(b_{it}\) when \(z_t=0\)
and \(\neg b_{it}\) when \(z_t=1\). The usual CNF encoding is
\begin{equation}\label{eq:bit-php}
 \BPHP_n^{n+1}
  =\bigwedge_{\substack{i<i'\in[m]\\z\in\{0,1\}^{\ell}}}
    \left(\bigvee_{t\in[\ell]}[b_{it}\ne z_t]
       \ \vee\ \bigvee_{t\in[\ell]}[b_{i't}\ne z_t]\right).
\end{equation}
It has \(v=m\ell\) variables and \(n\binom m2\) clauses of width \(2\ell\).
No separate range or pigeon axiom is required: every Boolean row encodes
exactly one available label.

\subsection{The proof system}
A \emph{linear clause} is a disjunction of affine equations over \(\F\).
The basic rules are complementary-parity resolution and semantic weakening:
\begin{equation}\label{eq:res-rules}
 \frac{A\vee(u=0)\qquad B\vee(u=1)}{A\vee B},
 \qquad
 \frac{C}{D}\quad\text{if } C\models D.
\end{equation}
The pivot \(u\) may be any affine form. A refutation is a finite acyclic
derivation from initial clauses to the empty clause. Its size is the number
of nodes, including initial nodes, in its proof DAG. Earlier nodes may be
used repeatedly, and no restriction is placed on proof depth or on repeated
use of related pivots.

We also consider the convention permitting every sound inference from at
most two linear-clause premises. Lemma~\ref{lem:affine-cover} reduces this
convention to \eqref{eq:res-rules} with constant node overhead.
Fresh unconstrained propositional variables, if used, may be set to
constants throughout a proof. Extension axioms defining new predicates
are not part of the system considered here.

\begin{theorem}[Main theorem, proposed proof]\label{thm:main}
For every fixed real \(K>0\), there is \(\ell_0(K)\) such that for every
\(\ell\ge\ell_0(K)\), every DAG-like \(\Res\) refutation of
\(\BPHP_{2^\ell}^{2^\ell+1}\) has more than \(2^{K\ell}\) nodes.
The conclusion holds for both rule conventions above.
\end{theorem}

Thus the minimum line count is \(n^{\omega(1)}\) along powers of two.
This also implies a superpolynomial lower bound on total proof-description
size. The input length in \eqref{eq:bit-php} is polynomial in \(n\), so
using input length instead of \(n\) gives the same qualitative conclusion.

\subsection{Prior results and the scope of the claim}
Efremenko, Garl\'ik, and Itsykson \cite{EGI25} prove a
\(2^{\Omega(n^{1/3}/\log n)}\) size bound for regular \(\Res\) refutations
of bit PHP. Their unrestricted DAG-like result is an affine-width bound.
Bhattacharya, Chattopadhyay, and Dvo\v{r}\'ak \cite{BCD24} exhibit formulas
that are exponentially hard for bottom-regular \(\Res\), although they
have short proofs in general ordinary resolution. This separation
demonstrates the importance of removing regularity.

Bhattacharya and Chattopadhyay \cite{BC25}, Byramji and Impagliazzo
\cite{BI25}, and the merged STOC 2026 paper \cite{BCBI26} obtain
exponential bounds in regimes allowing nearly quadratic proof depth.
Their bounds retain a proof-depth restriction. Itsykson, Podolskii,
and Shekhovtsov \cite{IPS26} explicitly distinguish this progress from
the general superpolynomial size problem; their unrestricted size
consequence is quadratic.

The statement proposed here removes both regularity and proof-depth
restrictions. The literature review found no preceding theorem with that
scope, but it is not an exhaustive certification of novelty.
The full \(\mathrm{AC}^0[p]\)-Frege lower-bound problem remains outside
the conclusion of Theorem~\ref{thm:main}.

\subsection{Proof organization}
The proof has three interfaces.
First, matching moments give a stable ordinary-degree filtration for
functional unary PHP, and a separate signed-direction construction detects
every nonzero row-linear bit polynomial.
Second, a restriction-dimension argument supplies a common polynomial for
all high-rank affine blocks, while low-rank blocks admit exact coefficient
packing. A single weighted substitution removes the whole affine family.
Third, each source proof clause is assigned a complete extension block;
local PC derivations simulate the rules without a proof-height cost.

Sections~\ref{sec:algebra}--\ref{sec:clause-simulation} prove these interfaces
in full. Section~\ref{sec:main-proof} closes the parameters.
Appendix~\ref{sec:topology} includes the precise homological input behind
the matching extension. The notebook is a provenance and navigation aid;
none of its lemmas is required as an unexpanded black box in this draft.
